% !TEX root = ../../../main.tex
\chapter{Sequences of real numbers}
\label{ch:analysis-i:sequences-of-real-numbers}

The foundations chapter introduced limits as a tool for completeness
and iteration. We now develop these ideas through sequence examples in more
detail. Limit laws and the squeeze principle come first; monotonicity
then gives convergence when the limit is not yet known. The later
sections recall the completeness results and tail bounds needed for series.

% !TEX root = ../../../main.tex
% Week 4, IMG_9226--IMG_9231. Definitions and examples are expanded from
% the handwritten notes; the limit-law proofs already appear in Chapter 1.
\section{Convergent sequences}
\label{sec:analysis-i:convergent-sequences}

A sequence is a list of real numbers indexed by the natural numbers.
Convergence describes what happens to its terms after sufficiently many
of them have been discarded. The first few terms do not decide its limit.

\begin{definition}
  A sequence $(a_n)$ \emph{converges} to $L\in\R$ if
  \[
    \forall\varepsilon>0\ \exists N\in\N\quad
    n\ge N\ \Longrightarrow\ |a_n-L|<\varepsilon.
  \]
  We write $a_n\to L$ or $\lim_{n\to\infty}a_n=L$.
\end{definition}

The index $N$ may depend on $\varepsilon$. Once it has been chosen,
\emph{every} term with index at least $N$ must lie in
$(L-\varepsilon,L+\varepsilon)$, as in
\autoref{fig:sequence-convergence-band}. Finding a few terms near $L$
would not be enough.

\begin{marginfigure}
  \centering
  % !TEX root = ../../main.tex
% Shared physical dimensions for interval and number-line diagrams.
\tikzset{
  interval diagram/.style={
    x=1cm,y=6mm,font=\footnotesize,
    color=black,line width=0.5pt,line cap=round,line join=round,
    >={Stealth[length=4pt,width=3pt]}
  },
  interval endpoint/.style={
    circle,draw=black,line width=0.5pt,
    minimum size=4pt,inner sep=0pt,outer sep=0pt
  },
  interval open/.style={interval endpoint,fill=white},
  interval closed/.style={interval endpoint,fill=black},
  interval label/.style={inner sep=1pt}
}

\begin{tikzpicture}[interval diagram,x=2.5mm,y=19mm]
  \fill[black!7] (4,0.7) rectangle (12,1.3);
  \draw[gray,dashed] (0,0.7) -- (12,0.7);
  \draw[gray,dashed] (0,1.3) -- (12,1.3);
  \draw[gray] (0,1) -- (12,1);
  \draw[gray,dashed] (4,0) -- (4,1.55);
  \draw[->] (0,0) -- (12.7,0) node[right=4pt] {$n$};
  \draw[->] (0,0) -- (0,1.8) node[above=4pt] {$a_n$};
  \node[left=5pt] at (0,0.7) {$L-\varepsilon$};
  \node[left=5pt] at (0,1) {$L$};
  \node[left=5pt] at (0,1.3) {$L+\varepsilon$};
  \node[below=6pt] at (4,0) {$N$};
  \foreach \n in {1,...,12}{
    \pgfmathsetmacro{\term}{1+(-1)^\n/\n}
    \fill (\n,\term) circle[radius=1.4pt];
  }
\end{tikzpicture}

  \caption[An epsilon band for a convergent sequence.]{For $a_n=1+(-1)^n/n$, all terms from $N=4$ onward lie in
    the band $1-0.3<a_n<1+0.3$. A smaller band may require a later $N$.}
  \label{fig:sequence-convergence-band}
\end{marginfigure}

\subsection{Rules for limits}

If $a_n\to L$, $b_n\to M$, and $c\in\R$, then
\begin{align*}
  ca_n&\to cL, & a_n\pm b_n&\to L\pm M,\\
  a_nb_n&\to LM, & \frac{a_n}{b_n}&\to\frac{L}{M}\quad(M\ne0).
\end{align*}
In the quotient rule the denominator is nonzero for all sufficiently
large $n$, because $b_n\to M\ne0$. These are the rules proved in
\autoref{thm:analysis-i:limit-laws}. The scalar rule follows from the
product rule by using the constant sequence $c$.

% IMG_9226 asks for direct proofs. Present the key estimates as a worked
% example, as requested by the author, rather than a classroom task.
\begin{example}[Bounds in the product and quotient rules]
  \label{ex:analysis-i:limit-law-bounds}
  Suppose $a_n\to L$ and $b_n\to M$. Eventually $|a_n-L|<1$, so
  $|a_n|\le A$ with $A=|L|+1$. The product error satisfies
  \[
    |a_nb_n-LM|\le A|b_n-M|+|M||a_n-L|.
  \]
  Given $\varepsilon>0$, choose $N$ so that this bound on $|a_n|$
  holds and the two errors are respectively less than
  $\varepsilon/(2A)$ and $\varepsilon/(2(|M|+1))$ for $n\ge N$.
  The right-hand side is then less than $\varepsilon$, proving
  $a_nb_n\to LM$ directly from the definition.

  If $M\ne0$, eventually $|b_n-M|<|M|/2$, which gives
  $|b_n|\ge |M|-|b_n-M|>|M|/2$. Hence
  \[
    \left|\frac{a_n}{b_n}-\frac{L}{M}\right|
    \le\frac{2}{|M|}|a_n-L|
      +\frac{2|L|}{|M|^2}|b_n-M|.
  \]
  Choose $N$ so that the denominator bound holds and the two errors
  are respectively less than $\varepsilon|M|/4$ and
  $\varepsilon|M|^2/(4(|L|+1))$ for $n\ge N$. The displayed
  bound is then less than $\varepsilon$, proving the quotient rule.
\end{example}

\begin{theorem}[Squeeze principle]
  \label{thm:analysis-i:squeeze}
  Suppose that $a_n\le b_n\le c_n$ for all sufficiently large $n$.
  If $a_n\to L$ and $c_n\to L$, then $b_n\to L$.
\end{theorem}

\begin{proof}
  Subtracting $L$ gives
  \[
    -|a_n-L|\le a_n-L\le b_n-L\le c_n-L\le |c_n-L|.
  \]
  Consequently $|b_n-L|\le |a_n-L|+|c_n-L|$. Given
  $\varepsilon>0$, choose an index after which the inequalities hold and
  both errors on the right are less than $\varepsilon/2$. Then
  $|b_n-L|<\varepsilon$ for every later term, which proves convergence.
\end{proof}
% IMG_9228 has crossed sign changes. The valid signed-error bounds above
% replace them. IMG_9227's sketch interchanges the upper/lower labels.

The middle sequence may oscillate. What matters is that its available
room shrinks to zero, as \autoref{fig:sequence-squeeze} shows.
In particular, if $0\le a_n\le b_n$ and $b_n\to0$, then $a_n\to0$.

\begin{marginfigure}
  \centering
  % !TEX root = ../../main.tex
% Shared physical dimensions for interval and number-line diagrams.
\tikzset{
  interval diagram/.style={
    x=1cm,y=6mm,font=\footnotesize,
    color=black,line width=0.5pt,line cap=round,line join=round,
    >={Stealth[length=4pt,width=3pt]}
  },
  interval endpoint/.style={
    circle,draw=black,line width=0.5pt,
    minimum size=4pt,inner sep=0pt,outer sep=0pt
  },
  interval open/.style={interval endpoint,fill=white},
  interval closed/.style={interval endpoint,fill=black},
  interval label/.style={inner sep=1pt}
}

\begin{tikzpicture}[interval diagram,x=2.6mm,y=19mm]
  \draw[->] (0,0) -- (12.7,0) node[right=4pt] {$n$};
  \draw[->] (0,0) -- (0,2.25) node[above=4pt] {$y$};
  \draw[gray,dashed] (0,1) -- (12,1);
  \node[left=5pt] at (0,1) {$L$};
  \draw[gray,domain=1:12,samples=80] plot (\x,{1+1/\x});
  \draw[gray,domain=1:12,samples=80] plot (\x,{1-1/\x});
  \node at (4.8,1.55) {$c_n$};
  \node at (4.8,0.45) {$a_n$};
  \foreach \n in {1,...,12}{
    \pgfmathsetmacro{\term}{1+(-1)^\n/(2*\n)}
    \fill (\n,\term) circle[radius=1.4pt];
  }
  \fill (4,1.95) circle[radius=1.4pt];
  \node[right=5pt] at (4,1.95) {$b_n$ terms};
\end{tikzpicture}

  \caption[An oscillating sequence squeezed between two bounds.]{The terms $b_n=1+(-1)^n/(2n)$ are trapped between
    $a_n=1-1/n$ and $c_n=1+1/n$. The gray curves guide the eye;
    the sequence terms are discrete. Both bounds tend to $1$.}
  \label{fig:sequence-squeeze}
\end{marginfigure}

\subsection{Divergence to infinity}

\begin{definition}
  We write $a_n\to+\infty$ if
  \[
    \forall M\in\R\ \exists N\in\N\quad n\ge N\ \Longrightarrow\ a_n>M.
  \]
  Similarly, $a_n\to-\infty$ means that for every $M\in\R$, all
  sufficiently late terms satisfy $a_n<M$.
\end{definition}

These sequences do not have a finite real limit. Divergence to
$+\infty$ means eventually remaining above every prescribed level;
being unbounded alone does not imply this.

\subsection{Geometric sequences}

For a positive number $a$, the behaviour of $a^n$ depends on whether
$a$ is below, equal to, or above $1$.
\[
  a^n\longrightarrow
  \begin{cases}
    0,&0<a<1,\\
    1,&a=1,\\
    +\infty,&a>1.
  \end{cases}
\]
The binomial identity provides a useful estimate. For $n\in\N$,
\[
  (u+v)^n=\sum_{k=0}^n\binom nk u^{n-k}v^k,
  \qquad \binom nk=\frac{n!}{k!(n-k)!}.
\]
If $h>0$, all terms in the expansion of $(1+h)^n$ are nonnegative,
so $(1+h)^n\ge1+nh$.

For $0<a<1$, write $1/a=1+h$ with $h>0$. Then
\[
  0<a^n=\frac{1}{(1+h)^n}\le\frac{1}{1+nh}\longrightarrow0.
\]
The squeeze principle gives the limit. If $a>1$, write $a=1+h$.
Now $a^n\ge1+nh$, which eventually exceeds every $M\in\R$ by the
Archimedean property. Thus $a^n\to+\infty$.

\subsection{Taking increasingly high roots}

\begin{proposition}
  \label{thm:analysis-i:root-sequence-limits}
  We have $n^{1/n}\to1$ and, for every fixed $a>0$, $a^{1/n}\to1$.
\end{proposition}

\begin{proof}
  Set $\delta_n=n^{1/n}-1\ge0$. For $n\ge2$, keeping just the
  quadratic term in the binomial expansion gives
  \[
    n=(1+\delta_n)^n\ge\binom n2\delta_n^2,
    \qquad 0\le\delta_n^2\le\frac{2}{n-1}.
  \]
  Given $\varepsilon>0$, the right-hand side is eventually smaller
  than $\varepsilon^2$, so $0\le\delta_n<\varepsilon$.
  Therefore $n^{1/n}\to1$.

  If $a\ge1$, choose $n>a$. Monotonicity of positive $n$th roots gives
  $1\le a^{1/n}\le n^{1/n}$, so the squeeze principle gives
  $a^{1/n}\to1$. If $0<a<1$, apply this result to $1/a>1$ and use
  the quotient rule to obtain
  \[
    a^{1/n}=\frac{1}{(1/a)^{1/n}}\longrightarrow1.
  \]
\end{proof}

\autoref{fig:root-sequences} plots actual terms of two root sequences.
The base in $n^{1/n}$ grows, yet taking the $n$th root brings the terms
back towards $1$.

\begin{figure*}[htbp]
  \centering
  % !TEX root = ../../main.tex
% Shared physical dimensions for interval and number-line diagrams.
\tikzset{
  interval diagram/.style={
    x=1cm,y=6mm,font=\footnotesize,
    color=black,line width=0.5pt,line cap=round,line join=round,
    >={Stealth[length=4pt,width=3pt]}
  },
  interval endpoint/.style={
    circle,draw=black,line width=0.5pt,
    minimum size=4pt,inner sep=0pt,outer sep=0pt
  },
  interval open/.style={interval endpoint,fill=white},
  interval closed/.style={interval endpoint,fill=black},
  interval label/.style={inner sep=1pt}
}

\begin{tikzpicture}[interval diagram,x=5.5mm,y=25mm]
  \draw[->] (0,0) -- (21.2,0) node[right=5pt] {$n$};
  \draw[->] (0,0) -- (0,2.2) node[above=5pt] {$y$};
  \draw[gray,dashed] (0,1) -- (20.5,1);
  \node[left=5pt] at (0,1) {$1$};
  \node[left=5pt] at (0,2) {$2$};
  \foreach \n in {1,10,20}{
    \draw (\n,0) -- +(0,-2pt);
    \node[below=6pt] at (\n,0) {$\n$};
  }
  \foreach \n in {1,...,20}{
    \pgfmathsetmacro{\growingroot}{\n^(1/\n)}
    \pgfmathsetmacro{\fixedroot}{2^(1/\n)}
    \fill (\n,\growingroot) circle[radius=1.6pt];
    \fill[gray] (\n,\fixedroot) circle[radius=1.6pt];
  }
  \fill (4,-0.5) circle[radius=1.6pt];
  \node[right=6pt] at (4,-0.5) {$n^{1/n}$};
  \fill[gray] (12,-0.5) circle[radius=1.6pt];
  \node[right=6pt] at (12,-0.5) {$2^{1/n}$};
\end{tikzpicture}

  \caption[Two root sequences approaching one.]{Terms of $n^{1/n}$ and $2^{1/n}$ on Cartesian axes.
    Both sequences approach the horizontal line $y=1$. Only integer
    indices represent sequence terms.}
  \label{fig:root-sequences}
\end{figure*}

\subsection{Exercises on root limits}

The following exercises combine root limits with geometric growth and
the squeeze principle.

% !TEX root = ../../hw04.tex
% Source: HW4 Intro to Math Analysis I.pdf, page 1, question 1.
\begin{exercise}
  \label{ex:hw04:polynomial-root-limit}
  \solutionlink{hw04:polynomial-root-limit}
  Find the limit
  \[
    \lim_{n\to\infty}(5n^3-2n+100)^{1/n}.
  \]
\end{exercise}

% !TEX root = ../../hw04.tex
% Source: HW4 Intro to Math Analysis I.pdf, page 1, question 2.
\begin{exercise}
  \label{ex:hw04:exponential-root-limit}
  \solutionlink{hw04:exponential-root-limit}
  Find the limit
  \[
    \lim_{n\to\infty}(a^n+n)^{1/n}
  \]
  for $0\le a\le1$ and $a>1$, respectively.
\end{exercise}


% !TEX root = ../../../main.tex
% Week 4, IMG_9232--IMG_9234. AM--GM and the identification with e are
% stated in the source; the two monotonicity arguments and bound are proved.
\section{Monotone bounded sequence theorem}
\label{sec:analysis-i:monotone-bounded-sequence-theorem}

\begin{theorem}[Monotone bounded sequence theorem]
  \label{thm:analysis-i:monotone-convergence}
  An increasing sequence bounded above converges to the supremum of
  its terms. A decreasing sequence bounded below converges to their infimum.
\end{theorem}

Here increasing and decreasing allow equality between consecutive terms.
The proof was given in \autoref{sec:analysis-i:completeness-equivalences}.
For example, if $L=\sup\{a_n\mid n\in\N\}$ and $(a_n)$ is increasing,
some $a_N$ exceeds $L-\varepsilon$. Every later term then satisfies
$L-\varepsilon<a_N\le a_n\le L$.

\subsection{The sequence \texorpdfstring{$(1+1/n)^n$}{(1+1/n) to the power n}}

Consider
\[
  c_n=\left(1+\frac1n\right)^n.
\]
We first establish a finite limit by proving monotonicity and an upper
bound. We then compare the binomial expansion with the factorial series
to identify the limit as $e$.

\begin{proposition}[Arithmetic--geometric mean inequality]
  \label{thm:analysis-i:am-gm}
  For an integer $k\ge1$ and positive real numbers $a_1,\ldots,a_k$,
  \[
    (a_1\cdots a_k)^{1/k}\le\frac{a_1+\cdots+a_k}{k}.
  \]
\end{proposition}

We use this inequality, abbreviated AM--GM, without proof.

\begin{example}[Convergence using AM--GM]
  \label{ex:analysis-i:binomial-sequence-am-gm}
  Apply \autoref{thm:analysis-i:am-gm} to $n+1$ positive numbers,
  one equal to $1$ and the other $n$ equal to $(n+1)/n$.
  Their arithmetic mean is $1+1/(n+1)$. Hence
  \[
    \left(\frac{n+1}{n}\right)^{n/(n+1)}
    \le1+\frac{1}{n+1}.
  \]
  Raising both sides to the power $n+1$ gives $c_n\le c_{n+1}$.
  Thus $(c_n)$ is increasing.

  We still need an upper bound. It is convenient to bound $c_{2n}$
  first. Put $b=1+1/(2n)$. Factoring the difference of powers gives
  \begin{align*}
    b^n-1
      &=(b-1)\sum_{j=0}^{n-1}b^j\\
      &\le n(b-1)b^{n-1}
       =\frac12b^{n-1}<\frac12b^n.
  \end{align*}
  The sum has $n$ terms, each at most $b^{n-1}$, including when $n=1$.
  Thus $b^n<2$ and $c_{2n}=(b^n)^2<4$. Since the sequence is
  increasing, $c_n\le c_{2n}<4$ for every $n$.
  \autoref{thm:analysis-i:monotone-convergence} now proves that $(c_n)$
  converges to a finite limit.
  % The first inequality must be non-strict when n=1 (IMG_9234).
\end{example}

Convergence alone does not determine the value of the limit. To find it,
we compare the binomial expansion with the factorial series defining $e$.

\begin{proposition}[The binomial limit]
  \label{thm:analysis-i:binomial-limit}
  \[
    \lim_{n\to\infty}\left(1+\frac1n\right)^n=e.
  \]
\end{proposition}
% Addition awaiting author review. The lecture establishes
% convergence but does not supply this identification with the earlier e.
\begin{pendingreview}
\begin{proof}[Identification of the limit]
  Let $L=\lim c_n$, whose existence has just been proved. The binomial
  expansion gives
  \[
    c_n=\sum_{k=0}^{n}\frac{1}{k!}
      \prod_{j=0}^{k-1}\left(1-\frac{j}{n}\right)
    \le\sum_{k=0}^{n}\frac{1}{k!}\le e.
  \]
  The product for $k=0$ is empty and equals $1$. Every other factor
  lies between $0$ and $1$. Taking limits gives $L\le e$, where $e$
  is the factorial-series sum defined in
  \autoref{sec:analysis-i:denseness}.

  Now fix an integer $m\ge0$. For every $n\ge\max\{m,1\}$, positivity
  of the remaining terms gives
  \[
    c_n\ge\sum_{k=0}^{m}\frac{1}{k!}
      \prod_{j=0}^{k-1}\left(1-\frac{j}{n}\right).
  \]
  This sum has a fixed number of terms. Each factor tends to $1$ as
  $n\to\infty$, so the finite product and sum rules give
  $L\ge\sum_{k=0}^{m}1/k!$. These inequalities hold for every $m$.
  The factorial partial sums tend to $e$, hence $L\ge e$.
  Together with $L\le e$, this proves $L=e$ without exchanging a
  limit with an infinite sum.
\end{proof}
\end{pendingreview}

There is also a direct proof of strict monotonicity using differences
of powers.

\begin{example}[Monotonicity by differences of powers]
  \label{ex:analysis-i:binomial-sequence-powers}
  If $0<a<b$ and $k\ge2$ is an integer, then
  \begin{align*}
    b^k-a^k
      &=(b-a)\bigl(b^{k-1}+b^{k-2}a+\cdots+a^{k-1}\bigr)\\
      &<k(b-a)b^{k-1}.
  \end{align*}
  Each summand in the brackets is at most $b^{k-1}$, and at least one
  is strictly smaller. Substitute
  \[
    a=1+\frac{1}{n+1},\qquad b=1+\frac1n,\qquad k=n+1.
  \]
  Since $k(b-a)=1/n$, we obtain
  \[
    \left(1+\frac1n\right)c_n-c_{n+1}<\frac{c_n}{n},
  \]
  and therefore $c_n<c_{n+1}$.
  % The strict power bound needs 0<a<b and k>=2, omitted in IMG_9233.
\end{example}

% The supplied week summary has exponent 1/n, whereas IMG_9232 has n.
\begin{remark}[Changing the exponent]
  The related expression with exponent $1/n$ has a different limit.
  Indeed,
  \[
    1\le\left(1+\frac1n\right)^{1/n}\le2^{1/n}\longrightarrow1,
  \]
  so it tends to $1$ by the squeeze principle.
\end{remark}

The two-variable case of AM--GM admits a direct proof that also allows
zero inputs and determines when equality holds.

% !TEX root = ../../hw04.tex
% Source: HW4 Intro to Math Analysis I.pdf, page 2, question 8.
\begin{exercise}
  \label{ex:hw04:two-variable-am-gm}
  \solutionlink{hw04:two-variable-am-gm}
  Let $a,b\ge0$. Prove the arithmetic mean-geometric mean inequality
  \[
    \sqrt{ab}\le\frac{a+b}{2}
  \]
  and deduce the necessary and sufficient condition when the inequality
  above is strict. (Hint: Consider $(\sqrt a-\sqrt b)^2$.)
\end{exercise}


% !TEX root = ../../../main.tex
% Recall of the completeness result already proved in Chapter 1.
\section{Bolzano--Weierstrass theorem}
\label{sec:analysis-i:bolzano-weierstrass-theorem}

\begin{theorem}[Bolzano--Weierstrass]
  \label{thm:analysis-i:bolzano-weierstrass}
  Every bounded real sequence has a convergent subsequence.
\end{theorem}

\begin{marginfigure}[22mm]
  \centering
  % !TEX root = ../../main.tex
% Shared physical dimensions for interval and number-line diagrams.
\tikzset{
  interval diagram/.style={
    x=1cm,y=6mm,font=\footnotesize,
    color=black,line width=0.5pt,line cap=round,line join=round,
    >={Stealth[length=4pt,width=3pt]}
  },
  interval endpoint/.style={
    circle,draw=black,line width=0.5pt,
    minimum size=4pt,inner sep=0pt,outer sep=0pt
  },
  interval open/.style={interval endpoint,fill=white},
  interval closed/.style={interval endpoint,fill=black},
  interval label/.style={inner sep=1pt}
}

\begin{tikzpicture}[interval diagram,x=8mm,y=10mm]
  % Successive halves retain infinitely many indexed sequence terms.
  \draw[gray,dashed] (1.72,0.5) -- (1.72,-3.3);
  \node[above=5pt] at (1.72,0.5) {$L$};

  \draw (0,0) -- (4,0);
  \node[interval closed] at (0,0) {};
  \node[interval closed] at (4,0) {};
  \node[left=6pt] at (0,0) {$I_1$};
  \fill (3.5,0) circle[radius=1.4pt];
  \node[above=5pt] at (3.5,0) {$a_{n_1}$};

  \draw (0,-1.4) -- (2,-1.4);
  \node[interval closed] at (0,-1.4) {};
  \node[interval closed] at (2,-1.4) {};
  \node[left=6pt] at (0,-1.4) {$I_2$};
  \fill (0.45,-1.4) circle[radius=1.4pt];
  \node[above=5pt] at (0.45,-1.4) {$a_{n_2}$};

  \draw (1,-2.8) -- (2,-2.8);
  \node[interval closed] at (1,-2.8) {};
  \node[interval closed] at (2,-2.8) {};
  \node[left=6pt] at (0,-2.8) {$I_3$};
  \fill (1.28,-2.8) circle[radius=1.4pt];
  \node[above left=5pt] at (1.28,-2.8) {$a_{n_3}$};
  \node at (1.4,-3.8) {$\vdots$};
\end{tikzpicture}

  \caption[Selecting a subsequence in nested intervals.]{Schematic selection in nested intervals. Each $I_k$
    contains infinitely many terms, so we can choose $a_{n_k}\in I_k$
    with $n_1<n_2<\cdots$. The shrinking intervals share the point $L$;
    the construction continues beyond the rows shown.}
  \label{fig:bolzano-weierstrass-selection}
\end{marginfigure}

A subsequence $(a_{n_k})$ uses strictly increasing indices
$n_1<n_2<\cdots$. It keeps some terms in their original order.
The nested-interval proof in
\autoref{sec:analysis-i:completeness-equivalences} repeatedly retains
a half-interval containing infinitely many terms, then chooses a term
in each retained interval with index larger than the previous one.
\autoref{fig:bolzano-weierstrass-selection} shows these choices.
If $L$ is the common point of the retained intervals $I_k$, then
\[
  |a_{n_k}-L|\le |I_k|\longrightarrow0,
\]
where $|I_k|$ denotes the length of $I_k$. Thus the shrinking intervals
force the chosen subsequence to converge to $L$.

The theorem guarantees a convergent subsequence, rather than convergence
of the original sequence. For instance, $(-1)^n$ is bounded and has
constant subsequences with limits $1$ and $-1$, although the whole
sequence does not converge. These subsequential limits are the limit
points introduced in \autoref{sec:analysis-i:denseness}.

% !TEX root = ../../../main.tex
% Week 4, IMG_9234, together with the bounded-sequence treatment in §1.9.
% Extended tail bounds are supporting notation for the ratio and root tests.
\section{Upper and lower limits (lim sup and lim inf)}
\label{sec:analysis-i:upper-and-lower-limits}

For a bounded real sequence $(a_n)$, define its tail bounds by
\[
  \alpha_n=\inf\{a_k\mid k\ge n\},\qquad
  \beta_n=\sup\{a_k\mid k\ge n\}.
\]
Removing early terms can only raise the infimum and lower the supremum.
Thus $(\alpha_n)$ is increasing, $(\beta_n)$ is decreasing, and
\[
  \liminf_{n\to\infty}a_n=\lim_{n\to\infty}\alpha_n,
  \qquad
  \limsup_{n\to\infty}a_n=\lim_{n\to\infty}\beta_n.
\]
By \autoref{thm:analysis-i:extreme-limit-points}, these are respectively
the smallest and largest limits of convergent subsequences. The full
sequence converges exactly when these two numbers agree. The tail
picture and its proof appear in \autoref{sec:analysis-i:denseness}.

\subsection{Allowing infinite tail bounds}

The tests for series also use sequences that need not be bounded.
We allow tail bounds to take values in $\R\cup\{-\infty,+\infty\}$.
An unbounded-above tail has supremum $+\infty$, and an
unbounded-below tail has infimum $-\infty$. The definitions become
\[
  \liminf a_n=\sup_n\alpha_n,\qquad
  \limsup a_n=\inf_n\beta_n.
\]
For bounded sequences these agree with the preceding definitions.
The infinite values record divergence of the tail bounds; they are
not additional real numbers.

\begin{marginfigure}
  \centering
  % !TEX root = ../../main.tex
% Shared physical dimensions for interval and number-line diagrams.
\tikzset{
  interval diagram/.style={
    x=1cm,y=6mm,font=\footnotesize,
    color=black,line width=0.5pt,line cap=round,line join=round,
    >={Stealth[length=4pt,width=3pt]}
  },
  interval endpoint/.style={
    circle,draw=black,line width=0.5pt,
    minimum size=4pt,inner sep=0pt,outer sep=0pt
  },
  interval open/.style={interval endpoint,fill=white},
  interval closed/.style={interval endpoint,fill=black},
  interval label/.style={inner sep=1pt}
}

\begin{tikzpicture}[interval diagram,x=3.8mm,y=2.6mm]
  \node at (4.5,9.8) {$a_n=n$};
  \draw[->] (0,0) -- (7.6,0) node[right=4pt] {$n$};
  \draw[->] (0,0) -- (0,8.1) node[above=5pt] {$a_n$};
  \fill[black!7] (3,2.5) rectangle (7.1,7.8);
  \draw[gray,dashed] (0,2.5) -- (7.1,2.5);
  \node[left=5pt] at (0,2.5) {$r$};
  \draw[gray,dashed] (3,0) -- (3,7.8);
  \node[below=6pt] at (3,0) {$N$};
  \foreach \n in {1,...,6}{
    \fill (\n,\n) circle[radius=1.4pt];
  }
  \node at (7.5,7.2) {$\cdots$};

  \begin{scope}[yshift=-33mm,x=3.8mm,y=1.5mm]
    \node at (4.5,11.4) {$b_n=(-1)^n n$};
    \draw[->] (0,0) -- (7.6,0) node[right=4pt] {$n$};
    \draw[->] (0,-7.5) -- (0,8.1) node[above=5pt] {$b_n$};
    \foreach \n in {1,...,6}{
      \pgfmathsetmacro{\term}{(-1)^\n*\n}
      \fill (\n,\term) circle[radius=1.4pt];
    }
    \node at (8,8.3) {$\cdots$};
    \node at (8,-8.3) {$\cdots$};
  \end{scope}
\end{tikzpicture}

  \caption[Two patterns of infinite tail bounds.]{For $a_n=n$, every sufficiently late term stays above a
    chosen level $r$. Here $r=5/2$ and $N=3$.
    For $b_n=(-1)^n n$, every tail reaches arbitrarily far in both
    directions. Dots are actual terms; ellipses indicate continuation.}
  \label{fig:infinite-tail-bounds}
\end{marginfigure}

\begin{example}[Two patterns of unbounded tails]
  \label{ex:analysis-i:infinite-tail-bounds}
  For $a_n=n$, the tail beginning at $N$ has $\alpha_N=N$ and
  $\beta_N=+\infty$. Hence both $\liminf a_n$ and $\limsup a_n$
  equal $+\infty$, and $a_n\to+\infty$.
  For $b_n=(-1)^n n$, every tail is unbounded above and below, so
  $\alpha_N=-\infty$ and $\beta_N=+\infty$ for every $N$.
  Its lower and upper limits are different, and it does not tend to
  either infinity. \autoref{fig:infinite-tail-bounds} compares the two
  patterns. Infinite tail bounds describe the whole tail, even though
  each individual term is a finite real number.
\end{example}

The feature we will need is an eventual bound. If $\limsup a_n<r$,
then $\beta_N<r$ for some $N$, so $a_n<r$ for every $n\ge N$.
Similarly, if $\liminf a_n>r$, then every sufficiently late term
exceeds $r$. The strict gap between the limit bound and $r$ is essential.

\begin{proposition}
  If $a_n\le b_n$ for every sufficiently large $n$, then
  $\liminf a_n\le\limsup b_n$.
\end{proposition}

\begin{proof}
  For all sufficiently late tails, $\alpha_n^{a}\le a_k\le b_k
  \le\beta_n^{b}$ whenever $k\ge n$. Given any two tail indices
  $i,j$, choose $k$ at least as large as both and beyond the index
  where $a_k\le b_k$. Then $\alpha_i^{a}\le\beta_j^{b}$.
  Taking the supremum over $i$ and then the infimum over $j$ gives
  the claimed inequality, including when a bound is infinite.
\end{proof}

% !TEX root = ../../../main.tex
\section{Cauchy sequences}
\label{sec:analysis-i:cauchy-sequences}

% BEGIN TUFTE PLACEHOLDER: supplied sample, lines 623--639.
\begingroup\sloppy
% Replace this entire specimen block with reviewed course notes.
References appear as bracketed numbers in the text, using the normal
\doccmddef{cite} command. Full entries are collected in the bibliography
before the appendices by \doccmddef{printbibliography}, in order of first
citation. The build runs Biber to resolve the entries in
\texttt{bibliography.bib}.

To enter multiple citations at one location~\cite{Tufte2006,Tufte1990},
provide a list of keys separated by commas,
as in \Verb|\cite{Tufte2006,Tufte1990}|.
An optional argument supplies a page or other locator, rather than a
vertical offset.
\begin{docspec}
  \doccmd{cite}[\docopt{postnote}]\{\docarg{bibkey1,bibkey2,\ldots}\}
\end{docspec}

\lipsum[7]
\par\endgroup
% END TUFTE PLACEHOLDER

